% ============================================================================
% CHAPTER 16: PROBLEM SET
% ============================================================================

\chapter{Problem Set}
\label{ch:problem_set}

This chapter consolidates all exercises from the modeling examples chapters. Problems are organized by domain to facilitate focused practice.


% ============================================================================

\begin{exercise}
A spring-mass system has $m = 2$ kg and $k = 200$ N/m. Find:
\begin{enumerate}[label=(\alph*)]
    \item The natural frequency in rad/s and Hz
    \item The transfer function $X(s)/F(s)$
    \item The response to a 10 N step force
\end{enumerate}
\end{exercise}

\begin{exercise}
Add a damper with $b = 8$ N$\cdot$s/m to the system in Exercise 1.
\begin{enumerate}[label=(\alph*)]
    \item Calculate the damping ratio
    \item Find the damped natural frequency
    \item Determine if the system is underdamped, critically damped, or overdamped
\end{enumerate}
\end{exercise}

\begin{exercise}
A motor with $J_m = 0.005$ kg$\cdot$m$^2$ drives a load with $J_L = 0.5$ kg$\cdot$m$^2$ through a 10:1 gear ratio.
\begin{enumerate}[label=(\alph*)]
    \item Calculate the reflected load inertia at the motor
    \item Find the total effective inertia seen by the motor
    \item Derive the transfer function from motor torque to load angle
\end{enumerate}
\end{exercise}

\begin{exercise}
For the quarter-car model, derive the transfer function from road velocity $\dot{X}_r(s)$ to body acceleration $s^2 X_s(s)$. This represents the acceleration felt by passengers.
\end{exercise}

\begin{exercise}
Linearize the simple pendulum equation about $\theta_0 = \pi/6$ (30 degrees). Find the equilibrium torque and the linearized dynamics.
\end{exercise}

\begin{exercise}
A tuned mass damper for building vibration control has the following parameters:
\begin{itemize}
    \item Building mass: $M = 1000$ kg
    \item Damper mass: $m = 20$ kg
    \item Damper spring constant: $k = 5000$ N/m
    \item Damper coefficient: $b = 500$ N$\cdot$s/m
\end{itemize}

\begin{enumerate}[label=(\alph*)]
    \item Calculate the natural frequency of the tuned damper
    \item If the building is excited at 2 Hz, find the phase lag of the damper response
    \item Design the spring constant to match the building frequency of 1.5 Hz
\end{enumerate}
\end{exercise}

\begin{exercise}
Consider a flexible robot arm modeled as a two-inertia system: motor inertia $J_m = 0.01$ kg$\cdot$m$^2$, link inertia $J_L = 0.5$ kg$\cdot$m$^2$, shaft torsional stiffness $\kappa = 50$ N$\cdot$m/rad, and bearing friction $B = 0.1$ N$\cdot$m$\cdot$s/rad.

\begin{enumerate}[label=(\alph*)]
    \item Derive the transfer function from motor torque to link angle
    \item Find the natural frequencies (radians and Hz)
    \item Explain why this system exhibits vibration and how to suppress it
\end{enumerate}
\end{exercise}

\begin{exercise}
\textbf{Concept:} Compare the responses of three second-order mechanical systems with the same natural frequency ($\omega_n = 10$ rad/s) but different damping ratios: $\zeta_1 = 0.3$, $\zeta_2 = 0.7$, $\zeta_3 = 1.0$.

\begin{enumerate}[label=(\alph*)]
    \item For each, describe the qualitative step response: overshoot, oscillations, settling time
    \item Rank them by settling time (2\% criterion)
    \item Which would you choose for a robotics application and why?
\end{enumerate}
\end{exercise}

\begin{exercise}
\textbf{Design:} A vibration isolator for precision equipment must satisfy:
\begin{itemize}
    \item Support mass: $m = 50$ kg
    \item Isolation efficiency $> 80\%$ at 10 Hz (reduce vibration amplitude by factor of 5)
    \item Static deflection $< 10$ cm at rated load
\end{itemize}

\begin{enumerate}[label=(\alph*)]
    \item Calculate the required spring constant
    \item Find the natural frequency
    \item Verify the isolation efficiency at 10 Hz (calculate transmissibility)
    \item Specify a damping ratio for practical implementation
\end{enumerate}
\end{exercise}

\begin{exercise}
\textbf{Concept:} The inverted pendulum on a cart is fundamentally unstable. Explain:

\begin{enumerate}[label=(\alph*)]
    \item Why feedback control is necessary (discuss the pole locations)
    \item What happens to the system response without control when perturbed
    \item Why a cart-based solution is simpler than a balance-wheel solution
\end{enumerate}
\end{exercise}

% ============================================================================

\begin{exercise}
For an RC circuit with $R = 47$ k$\Omega$ and $C = 0.22$ $\mu$F:
\begin{enumerate}[label=(\alph*)]
    \item Find the time constant and cutoff frequency
    \item Sketch the Bode plot (magnitude and phase)
    \item Find the output voltage 5 ms after a 10V step input
\end{enumerate}
\end{exercise}

\begin{exercise}
Design an RLC band-pass filter with center frequency 10 kHz and bandwidth 1 kHz. Calculate $R$, $L$, $C$, and $Q$.
\end{exercise}

\begin{exercise}
An op-amp integrator has $R = 100$ k$\Omega$ and $C = 1$ $\mu$F. If the input is $v_{\text{in}}(t) = 2\sin(100t)$ V:
\begin{enumerate}[label=(\alph*)]
    \item Find the transfer function
    \item Calculate the steady-state output amplitude and phase
    \item Write the time-domain output expression
\end{enumerate}
\end{exercise}

\begin{exercise}
For the series RLC circuit, find the current $i(t)$ response to a unit step voltage input when $R = 100$ $\Omega$, $L = 0.1$ H, and $C = 100$ $\mu$F. Is the circuit underdamped, critically damped, or overdamped?
\end{exercise}

\begin{exercise}
Design a second-order Sallen-Key low-pass filter with cutoff frequency 1 kHz and Butterworth response. Specify all component values.
\end{exercise}

\begin{exercise}
\textbf{Concept:} Compare the frequency response characteristics of three filters:
\begin{enumerate}[label=(\alph*)]
    \item A first-order low-pass filter with $\omega_c = 1000$ rad/s
    \item A second-order Butterworth low-pass filter with $\omega_c = 1000$ rad/s
    \item A first-order high-pass filter with $\omega_c = 1000$ rad/s
\end{enumerate}

For each, sketch the Bode magnitude plot and describe:
\begin{enumerate}[label=(\alph*)]
    \item The slope (dB/decade) in the stopband
    \item The phase shift at the cutoff frequency
    \item Which frequencies are attenuated and which are passed
\end{enumerate}
\end{exercise}

\begin{exercise}
A boost converter steps up a 12V input to 48V at a switching frequency of 200 kHz. Assume efficiency $\eta = 0.92$.

\begin{enumerate}[label=(\alph*)]
    \item Calculate the required duty cycle
    \item If the load draws 5A at 48V, find the input current
    \item Design the filter inductance to limit current ripple to $\pm 5\%$
\end{enumerate}
\end{exercise}

\begin{exercise}
\textbf{Design:} An active low-pass filter must attenuate 60 Hz power line noise by at least 40 dB while preserving signals below 50 Hz.

\begin{enumerate}[label=(\alph*)]
    \item Determine the minimum filter order required
    \item Choose between Butterworth, Chebyshev, and Bessel responses and justify your choice
    \item If using a second-order topology, calculate the component values for unity-gain implementation
    \item Estimate the phase lag at 20 Hz
\end{enumerate}
\end{exercise}

\begin{exercise}
\textbf{Computational:} For the series RLC circuit with $R = 10$ $\Omega$, $L = 1$ H, and $C = 1$ $\mu$F:

\begin{enumerate}[label=(\alph*)]
    \item Find the resonant frequency
    \item Calculate the quality factor $Q$
    \item Find the magnitude of impedance at resonance, at $\omega = \omega_n/2$, and at $2\omega_n$
    \item If driven by $v_{\text{in}}(t) = 100\sin(\omega_n t)$ V, find the steady-state current amplitude at resonance
\end{enumerate}
\end{exercise}

\begin{exercise}
\textbf{Concept:} The PID controller can be implemented using three op-amp stages (one proportional, one integrator, one differentiator).

\begin{enumerate}[label=(\alph*)]
    \item Draw a block diagram showing how three op-amp circuits combine to form a PID controller
    \item Explain the role of each stage in controlling a second-order system
    \item Discuss practical issues with the differentiator stage (e.g., noise) and propose solutions
\end{enumerate}
\end{exercise}

% ============================================================================

\begin{exercise}
A DC motor has $R_a = 1$ $\Omega$, $K_t = K_e = 0.05$ N$\cdot$m/A, $J = 2 \times 10^{-4}$ kg$\cdot$m$^2$, and $B = 10^{-4}$ N$\cdot$m$\cdot$s/rad. Neglect inductance.
\begin{enumerate}[label=(\alph*)]
    \item Derive the transfer function from voltage to velocity
    \item Find the steady-state speed for $V_a = 12$ V with no load
    \item Find the stall torque (torque at zero speed) for $V_a = 12$ V
\end{enumerate}
\end{exercise}

\begin{exercise}
For the motor in Exercise 1, a 5:1 gear reducer is added, and the load inertia is $J_L = 0.01$ kg$\cdot$m$^2$ with load damping $B_L = 0.01$ N$\cdot$m$\cdot$s/rad.
\begin{enumerate}[label=(\alph*)]
    \item Find the reflected load inertia and damping at the motor
    \item Derive the new transfer function from voltage to load velocity
    \item Compare the new time constant to the original
\end{enumerate}
\end{exercise}

\begin{exercise}
A voice coil actuator has $R = 5$ $\Omega$, $L = 1$ mH, $K_f = K_e = 10$ N/A, $m = 0.1$ kg, and $b = 2$ N$\cdot$s/m.
\begin{enumerate}[label=(\alph*)]
    \item Derive the transfer function from voltage to position
    \item Find all poles and determine if the system is stable
    \item Calculate the DC gain
\end{enumerate}
\end{exercise}

\begin{exercise}
A piezoelectric actuator has $d_{33} = 500$ pm/V, resonant frequency 30 kHz, and damping ratio 0.02.
\begin{enumerate}[label=(\alph*)]
    \item Write the transfer function
    \item Find the static displacement for 100V input
    \item Find the displacement amplitude at 20 kHz for 100V sinusoidal input
\end{enumerate}
\end{exercise}

\begin{exercise}
\textbf{Design:} A servo motor must position a robotic arm joint with:
\begin{itemize}
    \item Load inertia: $J_L = 0.02$ kg$\cdot$m$^2$
    \item Load damping: $B_L = 0.05$ N$\cdot$m$\cdot$s/rad
    \item Desired bandwidth: 20 rad/s
    \item Maximum allowed steady-state error: 0.01 rad
\end{itemize}

\begin{enumerate}[label=(\alph*)]
    \item Select a motor with appropriate torque and speed ratings
    \item Choose a gear ratio to match the bandwidth requirement
    \item Design a proportional feedback controller to meet the steady-state error specification
    \item Estimate the settling time (2\% criterion)
\end{enumerate}
\end{exercise}

\begin{exercise}
\textbf{Concept:} Explain the role of back-EMF in DC motor speed regulation.

\begin{enumerate}[label=(\alph*)]
    \item Derive how back-EMF creates a speed-dependent resistance to applied voltage
    \item Show that steady-state speed increases linearly with applied voltage
    \item Explain why increased load torque causes the motor to slow down
    \item Discuss how feedback control can maintain constant speed despite load variations
\end{enumerate}
\end{exercise}

\begin{exercise}
A stepper motor with 200 steps/revolution (1.8° per step) drives a load with moment of inertia $J = 10^{-4}$ kg$\cdot$m$^2$ and friction damping $B = 10^{-5}$ N$\cdot$m$\cdot$s/rad.

\begin{enumerate}[label=(\alph*)]
    \item Calculate the motor's step angle in radians
    \item If stepped at 1000 Hz, what is the steady-state rotational velocity?
    \item Calculate the acceleration at a 1000 rad/s$^2$ acceleration specification
    \item Explain the hunting and resonance issues in stepper motor control
\end{enumerate}
\end{exercise}

\begin{exercise}
\textbf{Computational:} For a brushless DC motor (BLDC) with field-oriented control:

\begin{enumerate}[label=(\alph*)]
    \item Derive the relationship between q-axis voltage and motor torque
    \item Show that the motor behaves like a first-order system in the current loop
    \item Calculate the current loop bandwidth for $L = 5$ mH, $R = 1$ $\Omega$, and proportional gain $K_p = 10$
    \item Find the settling time for a 10A step reference current
\end{enumerate}
\end{exercise}

\begin{exercise}
\textbf{Concept:} Compare direct drive motors versus geared motors for a positioning application.

\begin{enumerate}[label=(\alph*)]
    \item List advantages and disadvantages of each approach
    \item Show how gear ratio affects reflected inertia and bandwidth
    \item Explain the tradeoff between speed and torque
    \item In what applications would you choose each approach?
\end{enumerate}
\end{exercise}

% ============================================================================

\begin{exercise}
A metal block ($m = 5$ kg, $C_p = 500$ J/kg$\cdot$K) is heated by a 100 W heater and loses heat to 25°C ambient through convection ($hA = 5$ W/K).
\begin{enumerate}[label=(\alph*)]
    \item Write the differential equation for block temperature
    \item Find the thermal time constant
    \item Find the steady-state temperature
    \item How long to reach 90\% of steady-state from ambient?
\end{enumerate}
\end{exercise}

\begin{exercise}
\textbf{Thermal Resistance Network:} A system has two thermal resistances in series: $R_1 = 0.5$ K/W and $R_2 = 0.8$ K/W, with a heat source $\dot{Q} = 500$ W and ambient temperature $T_\infty = 20$°C. The total thermal capacitance is $C_{th} = 10$ kJ/K.
\begin{enumerate}[label=(\alph*)]
    \item Find the total thermal resistance
    \item Find the steady-state temperature rise above ambient
    \item If the system starts at ambient, find the temperature after 100 seconds
\end{enumerate}
\end{exercise}

\begin{exercise}
Two tanks in series each have $A = 1$ m$^2$. The first tank drains into the second through resistance $R_1 = 50$ s/m$^2$, and the second drains to atmosphere through $R_2 = 100$ s/m$^2$.
\begin{enumerate}[label=(\alph*)]
    \item Write the state-space model
    \item Find the transfer function from $Q_{\text{in}}$ to $h_2$
    \item Find the system poles and time constants
    \item Explain how the inter-tank interaction affects the response
\end{enumerate}
\end{exercise}

\begin{exercise}
\textbf{Tank Level Dynamics:} A conical tank has height-dependent area: $A(h) = \pi(rh/H)^2$ where $r = 0.5$ m is base radius and $H = 2$ m is total height. The outlet has resistance $R_f = 80$ s/m$^2$.
\begin{enumerate}[label=(\alph*)]
    \item Write the nonlinear ODE for level $h$
    \item For $h = 1$ m at steady-state with constant inflow, find $Q_{\text{in}}$
    \item Linearize the model about this operating point
\end{enumerate}
\end{exercise}

\begin{exercise}
A hydraulic servo system has a 10 cm diameter piston, 1000 kg load mass, and bulk modulus $\beta = 1.5 \times 10^9$ Pa. The chamber volume is 0.5 L.
\begin{enumerate}[label=(\alph*)]
    \item Calculate the hydraulic natural frequency
    \item If damping ratio is 0.3, find the damping coefficient
    \item Sketch the expected step response
\end{enumerate}
\end{exercise}

\begin{exercise}
\textbf{Heat Exchanger Thermal Dynamics:} A counter-flow heat exchanger receives hot fluid at 80°C (inlet) with flow rate 0.5 kg/s, and cold fluid at 20°C with flow rate 0.3 kg/s. The heat transfer coefficient is $UA = 2$ kW/K, and thermal capacitances are $C_h = C_c = 5$ kJ/K.
\begin{enumerate}[label=(\alph*)]
    \item Write the energy balance equations for both sides
    \item Find the steady-state outlet temperatures
    \item Determine the time constants of the hot and cold sides
\end{enumerate}
\end{exercise}

\begin{exercise}
A hot water tank ($V = 200$ L) receives water at 15°C at rate 0.1 L/s and has a 3 kW heater. Heat loss to ambient is $UA = 10$ W/K. Find the steady-state temperature and time constant. What is the response time to reach 95\% of steady-state?
\end{exercise}

% ============================================================================

\begin{exercise}
A CSTR with volume 500 L and flow rate 50 L/min processes a first-order reaction with $k = 0.1$ min$^{-1}$. Feed concentration is 2 mol/L.
\begin{enumerate}[label=(\alph*)]
    \item Find the steady-state exit concentration
    \item Find the residence time
    \item Write the transfer function from $C_{\text{in}}$ to $C$
    \item Find the time constant
\end{enumerate}
\end{exercise}

\begin{exercise}
A mixing tank blends two streams: Stream 1 at 100 L/min with 10\% salt, and Stream 2 at 50 L/min with 2\% salt. Tank volume is 1000 L.
\begin{enumerate}[label=(\alph*)]
    \item Find the steady-state exit concentration
    \item If Stream 1 concentration suddenly changes to 15\%, find the new steady-state and the time to reach 90\% of the change
\end{enumerate}
\end{exercise}

\begin{exercise}
An exothermic batch reactor starts at 25°C with 1000 mol of reactant. Heat of reaction is $-50$ kJ/mol, heat capacity is 4 kJ/kg$\cdot$K, mass is 1000 kg, and reaction rate constant is $k = 0.01$ min$^{-1}$ at 25°C with $E_a/R = 5000$ K.
\begin{enumerate}[label=(\alph*)]
    \item Write the coupled ODEs for concentration and temperature (adiabatic case)
    \item Estimate the adiabatic temperature rise if reaction goes to completion
\end{enumerate}
\end{exercise}

\begin{exercise}
\textbf{CSTR Reaction Kinetics:} A CSTR operates with a second-order reaction $2A \to B$ having rate law $r = kC_A^2$ with $k = 0.1$ L/(mol$\cdot$min). The reactor volume is 100 L and inlet concentration is 1 mol/L.
\begin{enumerate}[label=(\alph*)]
    \item For inlet flow 10 L/min, calculate the steady-state concentration
    \item If temperature increases causing $k$ to double, find the new steady-state
    \item Compare reaction rates at both operating points
\end{enumerate}
\end{exercise}

\begin{exercise}
\textbf{CSTR Temperature Control:} An exothermic reaction in a CSTR generates heat at rate $\dot{Q}_{\text{gen}} = (-\Delta H_r)k(T)C_A \cdot V$ where $-\Delta H_r = 50$ kJ/mol, $k(T)$ follows Arrhenius with $E_a/R = 8750$ K. The cooling jacket removes heat at $\dot{Q}_{\text{cool}} = UA(T - T_j)$ with $UA = 5$ kW/K.
\begin{enumerate}[label=(\alph*)]
    \item Write the energy balance ODE
    \item For $T_j = 350$ K (set at nominal operating point), find the steady-state temperature
    \item If $T_j$ suddenly drops to 300 K, predict the qualitative response
\end{enumerate}
\end{exercise}

\begin{exercise}
A process has FOPDT model with $K = 1.5$, $\tau = 120$ s, and $\theta = 30$ s.
\begin{enumerate}[label=(\alph*)]
    \item Write the transfer function
    \item Calculate $\theta/\tau$ and comment on controllability
    \item Use first-order Padé approximation and find the system poles
\end{enumerate}
\end{exercise}

\begin{exercise}
\textbf{pH Control Challenge:} A strong acid-strong base neutralization uses $[H^+] = \frac{C_a - C_b + \sqrt{(C_a - C_b)^2 + 4K_w}}{2}$ where $K_w = 10^{-14}$.
\begin{enumerate}[label=(\alph*)]
    \item With acid concentration 0.1 M flowing at 100 mL/min and base 0.1 M at 95 mL/min, find pH at steady-state
    \item Calculate the sensitivity $\frac{d(\text{pH})}{dC_b}$ at this point
    \item If base flow increases to 100.5 mL/min, estimate the new pH and explain the control difficulty
\end{enumerate}
\end{exercise}

% ============================================================================

\begin{exercise}
An aircraft has short-period approximation:
\begin{equation}
\frac{\alpha(s)}{\delta_e(s)} = \frac{-2}{s^2 + 1.5s + 4}
\end{equation}
\begin{enumerate}[label=(\alph*)]
    \item Find the natural frequency and damping ratio
    \item Is this response acceptable for a fighter aircraft?
    \item Sketch the step response
\end{enumerate}
\end{exercise}

\begin{exercise}
A rocket has pitch moment of inertia $I_{yy} = 10^6$ kg$\cdot$m$^2$, thrust $T = 10^6$ N, and gimbal arm $\ell_T = 20$ m.
\begin{enumerate}[label=(\alph*)]
    \item Find the transfer function from gimbal angle to pitch rate
    \item Design a proportional controller to achieve 1 rad/s$^2$ pitch acceleration per degree of gimbal
\end{enumerate}
\end{exercise}

\begin{exercise}
A satellite reaction wheel has inertia $I_w = 0.01$ kg$\cdot$m$^2$ and the satellite has inertia $I_s = 100$ kg$\cdot$m$^2$ about one axis.
\begin{enumerate}[label=(\alph*)]
    \item If the wheel accelerates at 100 rad/s$^2$, find the satellite angular acceleration
    \item Find the maximum satellite rate if the wheel is limited to 6000 rpm
\end{enumerate}
\end{exercise}

\begin{exercise}
A quadrotor has mass 1.5 kg, arm length 0.25 m, and moment of inertia $I_{xx} = I_{yy} = 0.02$ kg$\cdot$m$^2$.
\begin{enumerate}[label=(\alph*)]
    \item Find the hover thrust per motor
    \item Derive the transfer function from motor force differential to roll angle
    \item If motor response is modeled as first-order with $\tau = 0.05$ s, find the complete roll transfer function
\end{enumerate}
\end{exercise}

\begin{exercise}
A satellite has inertia $I = 50$ kg$\cdot$m$^2$ about the pitch axis. Two reaction wheels are available with inertias $I_w = 0.5$ kg$\cdot$m$^2$ each, located 90\degree apart about the body axes.
\begin{enumerate}[label=(\alph*)]
    \item Write the coupled equations for pitch-yaw control using both wheels
    \item If the pitch wheel accelerates at 200 rad/s$^2$, find the satellite pitch angular acceleration
    \item Design a control law to achieve zero final angular rate when the pitch wheel reaches 3000 rpm
\end{enumerate}
\end{exercise}

\begin{exercise}
An aircraft has longitudinal dynamics with short-period natural frequency $\omega_n = 1.2$ rad/s and damping $\zeta = 0.45$. The phugoid has $\omega_p = 0.08$ rad/s and $\zeta_p = 0.06$.
\begin{enumerate}[label=(\alph*)]
    \item Sketch the pole-zero map for the longitudinal system
    \item Calculate the time constants for each mode
    \item What type of controller (P, PI, or PD) would be best for elevator control? Why?
\end{enumerate}
\end{exercise}

\begin{exercise}
A rocket undergoes pitch control with gimbal actuation. The rocket parameters are: $I_{yy} = 2 \times 10^6$ kg$\cdot$m$^2$, thrust $T = 5 \times 10^6$ N, gimbal arm $\ell = 10$ m, and aerodynamic damping coefficient $M_q = -3 \times 10^5$ N$\cdot$m$\cdot$s/rad.
\begin{enumerate}[label=(\alph*)]
    \item Derive the transfer function from gimbal angle $\delta$ to pitch rate $q$
    \item Find the system poles and determine stability
    \item Design a proportional controller to limit pitch rate to 5 deg/s for 1\degree gimbal command
\end{enumerate}
\end{exercise}

\begin{exercise}
A quadrotor must perform a roll maneuver from level flight. It has mass $m = 2$ kg, arm length $\ell = 0.3$ m, and $I_{xx} = 0.025$ kg$\cdot$m$^2$.
\begin{enumerate}[label=(\alph*)]
    \item For a rolling maneuver command of 0.5 rad/s$^2$, calculate the required roll torque
    \item Find the force differential between motors: $\Delta F = F_1 - F_2$
    \item If available motor thrust is 10 N each, what is the maximum roll acceleration?
\end{enumerate}
\end{exercise}

\begin{exercise}
Consider a satellite attitude control system using a three-axis reaction wheel assembly. The satellite inertia matrix is diagonal: $I_{xx} = 80$, $I_{yy} = 100$, $I_{zz} = 120$ kg$\cdot$m$^2$.
\begin{enumerate}[label=(\alph*)]
    \item Write the decoupled single-axis control law for each axis: $\tau = I \ddot{\theta}$
    \item If maximum wheel torque is 0.2 N$\cdot$m, find the maximum angular acceleration for each axis
    \item Design a PD controller for the yaw axis: $\tau = K_p e + K_d \dot{e}$, with desired natural frequency 0.1 rad/s and damping ratio 0.7
\end{enumerate}
\end{exercise}

% ============================================================================

\begin{exercise}
A drug has one-compartment kinetics with $V = 35$ L and $t_{1/2} = 4$ hours.
\begin{enumerate}[label=(\alph*)]
    \item Calculate the elimination rate constant
    \item Find the loading dose to achieve 20 mg/L immediately
    \item Find the maintenance infusion rate
    \item Write the transfer function from infusion rate to concentration
\end{enumerate}
\end{exercise}

\begin{exercise}
A patient's arterial compliance is $C = 2$ mL/mmHg and peripheral resistance is $R = 1$ mmHg$\cdot$s/mL.
\begin{enumerate}[label=(\alph*)]
    \item Find the time constant of the arterial system
    \item If cardiac output steps from 80 to 100 mL/s, find the new steady-state pressure
    \item Sketch the pressure response to the step change
\end{enumerate}
\end{exercise}

\begin{exercise}
A mechanical ventilator delivers pressure to a patient with $R_{aw} = 10$ cmH$_2$O$\cdot$s/L and $C_L = 0.05$ L/cmH$_2$O.
\begin{enumerate}[label=(\alph*)]
    \item Find the time constant
    \item If inspiratory pressure is 20 cmH$_2$O above PEEP, find the tidal volume
    \item Find the inspiratory time needed to achieve 95\% of the tidal volume
\end{enumerate}
\end{exercise}

\begin{exercise}
An artificial pancreas controller must regulate glucose. The simplified model has:
\begin{equation}
\frac{G(s)}{U(s)} = \frac{-2}{(10s+1)(50s+1)}
\end{equation}
where $G$ is glucose (mg/dL) and $U$ is insulin (U/hr), with time in minutes.
\begin{enumerate}[label=(\alph*)]
    \item Find the steady-state glucose change for 1 U/hr insulin increase
    \item Find the system poles and time constants
    \item Sketch the glucose response to a step insulin input
\end{enumerate}
\end{exercise}

\begin{exercise}
A two-compartment drug delivery system has: $V_1 = 10$ L, $V_2 = 40$ L, $k_{10} = 0.2$ hr$^{-1}$, $k_{12} = 0.15$ hr$^{-1}$, and $k_{21} = 0.1$ hr$^{-1}$.
\begin{enumerate}[label=(\alph*)]
    \item Calculate the drug concentration in both compartments at 4 hours following a 500 mg IV bolus dose
    \item Find the half-lives for the distribution and elimination phases
    \item If the drug is continuously infused at 10 mg/hr, find the steady-state concentrations in both compartments
\end{enumerate}
\end{exercise}

\begin{exercise}
A patient receiving mechanical ventilation has measured respiratory mechanics: $R_{aw} = 8$ cmH$_2$O$\cdot$s/L, $C_L = 0.04$ L/cmH$_2$O, $I_r = 0.03$ L/cmH$_2$O (chest wall compliance).
\begin{enumerate}[label=(\alph*)]
    \item Find the total respiratory system compliance and time constant
    \item For a constant inspiratory pressure of 15 cmH$_2$O (above PEEP), calculate tidal volume after 5 time constants
    \item Design a pressure control strategy to deliver 0.5 L tidal volume in 1 second
\end{enumerate}
\end{exercise}

\begin{exercise}
A pacemaker rate-response control system regulates heart rate based on activity. The transfer function from activity sensor voltage to pacing rate is:
\begin{equation}
\frac{R(s)}{A(s)} = \frac{K}{0.5s + 1}, \quad K = 100 \text{ bpm/V}
\end{equation}
\begin{enumerate}[label=(\alph*)]
    \item Calculate the steady-state pacing rate for activity input of 0.5 V
    \item Find the time to reach 90\% of steady-state rate
    \item Design a controller to achieve 0 to 120 bpm range with sensor output of 0 to 1 V
\end{enumerate}
\end{exercise}

\begin{exercise}
An infusion pump delivers IV medications with first-order actuator dynamics: $\frac{Q(s)}{V_{\text{cmd}}(s)} = \frac{K_p}{\tau_p s + 1}$, where $K_p = 5$ mL/min/V and $\tau_p = 0.1$ s.
\begin{enumerate}[label=(\alph*)]
    \item For a step command of 2 V, find the steady-state flow rate
    \item Calculate the actual flow rate at time $t = 0.2$ s
    \item If a patient needs 50 mL of drug infused in exactly 10 minutes, determine the required command voltage accounting for actuator dynamics
\end{enumerate}
\end{exercise}

\begin{exercise}
A prosthetic limb controller uses proportional myoelectric control, with muscle EMG signal driving joint angle:
\begin{equation}
\theta(s) = \frac{G(s)}{1 + G(s)H(s)}E(s)
\end{equation}
where $G(s) = \frac{K}{\tau s + 1}$ (actuator), $H(s) = K_f s$ (feedback with damping), $K = 45$ deg/V, $\tau = 0.2$ s, and $K_f = 0.1$ V$\cdot$s/deg.
\begin{enumerate}[label=(\alph*)]
    \item Find the closed-loop transfer function $\theta(s)/E(s)$
    \item For a step EMG input of 0.5 V, calculate steady-state joint angle
    \item Sketch the response and identify if overshoot or oscillation occurs
    \item Design a proportional controller to add to $G(s)$ that achieves 1 deg/V steady-state sensitivity with damping ratio 0.8
\end{enumerate}
\end{exercise}

\begin{exercise}
A glucose-insulin model for artificial pancreas design is based on the Bergman minimal model. Patient parameters are: $p_1 = 0.03$ min$^{-1}$, $p_2 = 0.02$ min$^{-1}$, $p_3 = 1.2 \times 10^{-5}$ mL$/$($\mu$U$\cdot$min$)$, $n = 0.1$ min$^{-1}$, $\gamma = 0.5$ mg$/$dL$/min$/($\mu$U$/mL$)$, with $G_b = 95$ mg/dL and $I_b = 12$ $\mu$U/mL.
\begin{enumerate}[label=(\alph*)]
    \item Linearize the system about the basal operating point $(G_0, I_0) = (G_b, I_b)$
    \item Find the transfer functions $G(s)/U(s)$ and $G(s)/D(s)$
    \item Design a proportional insulin infusion controller: $u(t) = K_p(G_{\text{ref}} - G(t))$ that regulates glucose to 100 mg/dL with settling time under 3 hours
\end{enumerate}
\end{exercise}

% ============================================================================

\begin{exercise}
A conveyor belt has equivalent inertia $J = 2$ kg$\cdot$m$^2$, friction $B = 0.5$ N$\cdot$m$\cdot$s/rad, drive roller radius $r = 0.1$ m, and motor constant $K_t = 0.5$ N$\cdot$m/A.
\begin{enumerate}[label=(\alph*)]
    \item Find the transfer function from motor current to belt velocity
    \item Calculate the current needed for 0.5 m/s belt speed
    \item Find the time constant
\end{enumerate}
\end{exercise}

\begin{exercise}
A single-link robot arm has $J = 0.5$ kg$\cdot$m$^2$, $B = 0.1$ N$\cdot$m$\cdot$s/rad, $m = 2$ kg, and $\ell = 0.5$ m.
\begin{enumerate}[label=(\alph*)]
    \item Write the nonlinear equation of motion
    \item Linearize about $\theta_0 = 0$ (horizontal)
    \item Find the gravity compensation torque
    \item Derive the transfer function from torque to angle (with gravity compensated)
\end{enumerate}
\end{exercise}

\begin{exercise}
An M/M/1 queue has arrival rate $\lambda = 80$ jobs/s and service rate $\mu = 100$ jobs/s.
\begin{enumerate}[label=(\alph*)]
    \item Calculate the utilization
    \item Find the mean queue length
    \item Find the mean response time (waiting + service)
    \item If the arrival rate increases to 90 jobs/s, how do the metrics change?
\end{enumerate}
\end{exercise}

\begin{exercise}
A data center has thermal capacitance $C = 10^6$ J/K, heat load $Q_{\text{IT}} = 100$ kW, and air flow rate that removes heat according to $\dot{Q}_{\text{cool}} = 5000(T - 18)$ W where $T$ is in °C.
\begin{enumerate}[label=(\alph*)]
    \item Find the steady-state temperature
    \item Calculate the thermal time constant
    \item If IT load suddenly increases to 120 kW, find the new steady-state and time to reach it
\end{enumerate}
\end{exercise}

\begin{exercise}
A traffic intersection has arrival rate 600 vehicles/hour and saturation flow 1800 vehicles/hour. The cycle length is 90 seconds.
\begin{enumerate}[label=(\alph*)]
    \item Find the minimum green time to prevent queue buildup
    \item If green time is 40 seconds, calculate the average queue at the start of green
    \item Model the queue dynamics as a hybrid system
\end{enumerate}
\end{exercise}

\begin{exercise}[CNC Positioning Control]
A CNC machine has proportional position controller gain $K_p = 50$ V/mm and velocity controller gain $K_i = 10$ A/(mm/s). The motor has torque constant $K_t = 0.8$ N$\cdot$m/A, axis inertia $J = 0.2$ kg$\cdot$m$^2$, and viscous damping $B = 2$ N$\cdot$m$\cdot$s/rad. The lead screw couples the motor to the axis with a 5 mm pitch.
\begin{enumerate}[label=(\alph*)]
    \item Convert the velocity command from the position controller into screw rotational velocity (rad/s)
    \item Write the motor dynamics in state-space form with velocity feedback
    \item Find the steady-state position error for a constant velocity ramp command of 10 mm/s
    \item Design the proportional gain $K_p$ to achieve a settling time of 0.5 s for a 1 mm step command
    \item Explain how cutting force disturbance $F_{\text{cut}} = 500$ N affects steady-state accuracy
\end{enumerate}
\end{exercise}

\begin{exercise}[Queueing System Design]
A server farm handles web requests with M/M/1 queue dynamics. Requests arrive at $\lambda = 50$ req/s and each server processes requests at $\mu = 60$ req/s.
\begin{enumerate}[label=(\alph*)]
    \item Calculate the server utilization and assess stability
    \item Compute the average number of requests in the system
    \item Find the average response time (waiting + processing)
    \item If we want to maintain $W < 50$ ms, how many additional servers are needed?
    \item Model the queue length as a first-order system: $\dot{q} = \lambda - \mu(t)$, and find the characteristic time constant
\end{enumerate}
\end{exercise}

\begin{exercise}[Data Center Thermal Control]
A data center has thermal dynamics:
\begin{equation}
C\frac{\dd T}{\dd t} = Q_{\text{IT}} - Q_{\text{cool}}(T, \dot{m})
\end{equation}

where $C = 5 \times 10^6$ J/K is the room thermal capacitance, $Q_{\text{IT}} = 500$ kW is the IT heat load, and cooling is $Q_{\text{cool}} = \dot{m} c_p (T - T_{\text{supply}})$ with $c_p = 1005$ J/(kg$\cdot$K), $T_{\text{supply}} = 15°$C, and controlled mass flow $\dot{m}(t)$ in kg/s.
\begin{enumerate}[label=(\alph*)]
    \item Find the steady-state room temperature when $\dot{m} = 100$ kg/s
    \item Linearize the system about this operating point
    \item Design a proportional controller for mass flow: $\dot{m}(t) = K_p(T_{\text{ref}} - T(t))$ with $T_{\text{ref}} = 20°$C
    \item Find the time constant of the closed-loop system with $K_p = 5$ kg/(s$\cdot$°C)
    \item Analyze stability: what maximum $K_p$ prevents oscillation (critical damping)?
\end{enumerate}
\end{exercise}

\begin{exercise}[Network Congestion Window Dynamics]
TCP congestion control uses a window-based mechanism. The congestion window $W$ (in packets) evolves according to:
\begin{equation}
\frac{\dd W}{\dd t} = \frac{A}{W} - B \cdot p(W)
\end{equation}

where $A$ and $B$ are constants, and $p(W)$ is the packet loss probability, modeled as $p(W) = 0.3 \cdot (W/100)^2$ when $W > 50$ and $p(W) = 0$ otherwise (no loss during congestion avoidance).
\begin{enumerate}[label=(\alph*)]
    \item Find the equilibrium window size when $A = 2$, $B = 1$
    \item Analyze stability of the equilibrium: is it stable or unstable?
    \item If the RTT suddenly increases by 20\%, how does the equilibrium window change?
    \item Model a step increase in competing traffic (increase in $B$) and sketch the expected transient response
    \item Propose a controller to maintain target window size $W_{\text{ref}} = 80$ packets
\end{enumerate}
\end{exercise}

\begin{exercise}[Two-Layer Queueing Network]
Consider a tandem queueing network where jobs arrive at Queue 1 (servers rate $\mu_1 = 100$ jobs/s), then proceed to Queue 2 (servers rate $\mu_2 = 80$ jobs/s). The arrival rate is $\lambda = 70$ jobs/s.
\begin{enumerate}[label=(\alph*)]
    \item Are both queues stable? Justify your answer
    \item Calculate the average delay at Queue 1 and Queue 2 separately
    \item Find the total average time a job spends in the network
    \item If Queue 2 becomes a bottleneck ($\mu_2$ reduced to 65 jobs/s), what is the new stability boundary?
    \item Propose a feedback control strategy (admission control) to maintain both queues stable despite temporary traffic spikes
    \item Sketch a block diagram showing fluid-flow models of both queues with control feedback
\end{enumerate}
\end{exercise}
